% ===================================================================
%  TABELO N° 5 — La Domo e lua Konstrukto.
%  Feuillets 32 a 38 du fac-simile = folios 30 a 36.
%  Correspondance : folio imprime = numero de feuillet - 2.
%
%  Ca tabelo esas mayoritatim un DIALOGO inter Ioannes e lua onklulo
%  (patre lu artikitekto, siorulo Blanko l'entraprezisto, siorulo Rufo
%  l'arkitekto). Nula alinea esas numerizita da la autoro : la unika
%  numeroro qua aparas en la paje esas ta di la kayerar-signaturo
%  (« 3. Ido. » ye la pedo di la folio 33), qua ne aparteni a la texto
%  e ne esas transkriptita. Onu do klavizas la tota tabelo sub la
%  numero d'alinea « 01 », singla replika o alinea sequanta konstitu-
%  ante un rango sequanta, quale la konsigno preskriptas por alineati
%  ne-numerizita.
% ===================================================================

% --- Feuillet 32 (folio 30, ouverture de tableau) -------------------
\begin{VUpage}[32]{30}
%%K t05-apar-1 apar
\VUsaut{6.0mm}
\VUtitre{15.0pt}{60}{TABELO N\textsuperscript{o} 5}
\VUsaut{1.4mm}
\VUtitre{11.4pt}{0}{La Domo e lua Konstrukto.}
\VUsaut{2.2mm}

%%K t05-tit-1 sub
\VUpk{10.2pt}{40}{Bona renkontro.}
\VUsaut{3.0mm}

%%K t05-01-1 p
\VUblancAlinea
\textsc{Ioannes}. --- Onklo mea, me tre deziras savar\nl
quale on konstruktas \VUgras{domo}.

%%K t05-01-2 p
\VUblancAlinea
\textsc{La Onklulo} (80). --- To esas tre facila, mea\nl
amiko, venez kun me, me montros a tu ta\nl
quan sioro Bernardus (79) igas konstruktar e\nl
quan me habitos.

%%K t05-01-3 p
\VUblancAlinea
\textsc{Ioannes}. --- Qua do desegnis la \VUgras{plano} \textsuperscript{(76)} di\nl
ca domo?

%%K t05-01-4 p
\VUblancAlinea
\textsc{L'Onklulo}. --- \VUgras{L'arkitekto} \textsuperscript{(75)}; il prizentis\nl
tre klara desegnuro qua esis aprobata, e \VUgras{l'en}\cc
\VUgras{traprezisto} \textsuperscript{(77)} komencis la labori.

%%K t05-01-5 p
\VUblancAlinea
\textsc{Ioannes}. --- Qua do esas l'entraprezisto?

%%K t05-01-6 p
\VUblancAlinea
\textsc{L'Onklulo}. --- Lu esas siorulo Blanko, la\nl
patrulo di tua amiko Iakobus. Il direktas la\nl
laboristi kun siorulo Rufo, l'arkitekto.

%%K t05-01-7 p
\VUblancAlinea
Nu! Yen justatempe venas siorulo Blanko\nl
kun sua \VUgras{linealo} \textsuperscript{(78)}, e siorulo Rufo kun sua\nl
plano.

%%K t05-01-8 p
\VUblancAlinea
Bon jorno, siori. Quale vi standas?

%%K t05-01-9 p
\VUblancAlinea
\textsc{Siorulo Rufo}. --- Tre bone, danko. Bon\nl
jorno, Ioannes, quante kreskis ica infanto!

%%K t05-01-10 p
\VUblancAlinea
\textsc{L'Onklulo}. --- Yes, vere, il esas mikra viro.\nl
Il esas tre serioza, on mustas informar lu pri\nl
omno quon il vidas. Cadie, il volus savar\nl
quale on konstruktas domo.
\end{VUpage}

% --- Feuillet 33 (folio 31) ------------------------------------------
\begin{VUpage}[33]{31}
%%K t05-tit-2 sub
\VUpk{10.2pt}{40}{La laboristi.}
\VUsaut{3.0mm}

%%K t05-01-11 p
\VUblancAlinea
\textsc{Siorulo Blanko}. --- Nu, venez kun me, ami\cc
keto, me quik explikos lo a tu, nam me tre\nl
amas la pueri qui volas instruktar su.

%%K t05-01-12 p
\VUblancAlinea
Tu vidas ta laboristo qua tenas \VUgras{shovelo} \textsuperscript{(82)}\nl
en manuo: il esas \VUgras{ter-laboristo} \textsuperscript{(81)}. Il exka\cc
vas \VUgras{trancheo} \textsuperscript{(46)} por pozar la aquo-tubaro. Il\nl
anke exkavis la sulo en la loko ube esas la\nl
domo, por ke la masonisto erektez la quar\nl
\VUgras{muri}. La parto di ta muri qua esas en la\nl
sulo esas nomizata \VUgras{fundamenti} \textsuperscript{(2)}. La spaco\nl
inkluzita inter la fundamenti formacas la\nl
\VUgras{kelero} \textsuperscript{(3)}. Ta kelero havas mikra fenestro\nl
nomizata \VUgras{keler-fenestro} \textsuperscript{(5)}, \VUgras{pordo} \textsuperscript{(4)} ed eska\cc
lero.

%%K t05-01-13 p
\VUblancAlinea
La \VUgras{masonisto} \textsuperscript{(108)} duris konstruktar la muri\nl
lasante aperturi por la \VUgras{pordi} \textsuperscript{(8)} e la \VUgras{fenestri} \textsuperscript{(17)}.

%%K t05-01-14 p
\VUblancAlinea
\textsc{Ioannes}. --- Ma, ta masoniston me ne vidas!

%%K t05-01-15 p
\VUblancAlinea
S\textsuperscript{o} \textsc{Blanko}. --- Nun il finis laborar en la\nl
domo. Il esas proxim la \VUgras{kavaleyo} \textsuperscript{(54)}, sur sua\nl
\VUgras{eshafodo} \textsuperscript{(110)}, il konstruktas la muro di la \VUgras{lesi}\cc
\VUgras{veyo} \textsuperscript{(56)}.

%%K t05-01-16 p
\VUblancAlinea
Il havas trulo, trogo ed \VUgras{aplombo} \textsuperscript{(109)}. Ma tu\nl
ne memoros ta nomi.

%%K t05-01-17 p
\VUblancAlinea
\textsc{Ioannes}. --- Certe yes, nam me quik deman\cc
dos de mea onklulo mikra trulo e trogo, e me\nl
ipsa fabrikos aplombo per kordeto.

%%K t05-tit-3 sub
\VUsaut{3.0mm}
\VUpk{10.2pt}{40}{La materialo.}
\VUsaut{3.0mm}

%%K t05-01-18 p
\VUblancAlinea
\textsc{L'Onklulo}. --- Ha! tre bone. Ma dicez a me,\nl
Ioannes, quon on bezonas por konstruktar\nl
domo?
\end{VUpage}

% --- Feuillet 34 (folio 32) ------------------------------------------
\begin{VUpage}[34]{32}
%%K t05-01-19 p
\VUblancAlinea
\textsc{Ioannes}. --- On bezonas \VUgras{quadri} \textsuperscript{(59)} e \VUgras{mor}\cc
\VUgras{tero} \textsuperscript{(65)}.

%%K t05-01-20 p
\VUblancAlinea
\textsc{L'Onklulo}. --- Juste, videz ta viro, kun\nl
granda chapelo, sur sua \VUgras{charioto} \textsuperscript{(136)}. Il esas\nl
\VUgras{char-duktisto} \textsuperscript{(135)}. Singladie il adduktas hike\nl
\VUgras{petra bloki} \textsuperscript{(60)}. Il descharjas sua veturo en la\nl
korto, e la \VUgras{quadro-taliisti} \textsuperscript{(83)} talias la bloki\nl
e segas li per sua \VUgras{petro-segilo} \textsuperscript{(85)}. \VUgras{Manula}\cc
\VUgras{boristo} \textsuperscript{(146)} portas oli al masonisto qua lokizas\nl
li.

%%K t05-01-21 p
\VUblancAlinea
Dume, la \VUgras{dilutisto} \textsuperscript{(87)} preparas la mortero.\nl
Il bezonas \VUgras{sablo} \textsuperscript{(62)}, \VUgras{kalko} \textsuperscript{(63)} ed \VUgras{aquo} \textsuperscript{(64)}.\nl
La kalko esas proxim ilu en \VUgras{sako} \textsuperscript{(71)}, mason-\nl
\VUgras{servisto} \textsuperscript{(111)} adportas ad ilu la sablo sur\nl
\VUgras{brueto} \textsuperscript{(112)} e l'\VUgras{aprentiso} \textsuperscript{(113)} esas an la pum\cc
pilo (58). Il plenigas \VUgras{sitelo} \textsuperscript{(114)}. La mortero esas\nl
balde pronta. La \VUgras{morter-portisto} \textsuperscript{(107)} prenas lu\nl
ed adsupre portas, per la skalo, sur l'eshafodo,\nl
ube laboras masonisto qua finas ta latero di la\nl
domo.

%%K t05-01-22 p
\VUblancAlinea
E nun, quale esas nomizata la laboristo qua\nl
laboras la \VUgras{tekto} \textsuperscript{(31)}?

%%K t05-01-23 p
\VUblancAlinea
\textsc{Ioannes}. --- Il esas \VUgras{tektisto} \textsuperscript{(118)}.

%%K t05-tit-4 sub
\VUsaut{3.0mm}
\VUpk{10.2pt}{40}{La tekto di la domo.}
\VUsaut{3.0mm}

%%K t05-01-24 p
\VUblancAlinea
S\textsuperscript{o} \textsc{Blanko}. --- Yes, ma unesme venas la\nl
\VUgras{karpentisto} \textsuperscript{(115)}.

%%K t05-01-25 p
\VUblancAlinea
La karpentisto laboras la \VUgras{karpenturo} \textsuperscript{(35)}. Il\nl
asemblas la \VUgras{trabi} \textsuperscript{(37)} per \VUgras{stifti} \textsuperscript{(117)}. Ta labo\cc
risti, ibe supre, kun sua larja brachi velura,\nl
esas karpentisti. L'una tenas trabeto, e l'altra\nl
sekas stifto per sua \VUgras{hakilo} \textsuperscript{(116)}. Ta qua esas\nl
proxim la \VUgras{kameno} \textsuperscript{(41)} esas la tektisto. Il klo\cc
\parplein
\end{VUpage}

% --- Feuillet 35 (folio 33) ------------------------------------------
\begin{VUpage}[35]{33}
%%K t05-noto-1 noto
\VUnotes{143.2mm}{%
(*) \VUgras{Balk-o} = D. \textit{Balken}, E. \textit{joist}, F. \textit{solive}, I. \textit{travicello},\nl
R. \textit{balka}, S. Port. \textit{viga}.\nl
Teknikala radiko til nun ne adoptita, ma propozenda por\nl
tradukar la senco dil vorto F. \textit{solive}, qua ne esas trabo\nl
F. \textit{poutre}, nek trabeto. Ol esas ligna stango quadratigita\nl
qua suportas planko-sulo e plafono.%
}

%%K t05-01-25 p suite
\VUcontinue
vagas lati sur la (*) balki, ed \VUgras{ardezi} \textsuperscript{(66)} sur la\nl
lati. Kelkafoye il pozas teguli.

%%K t05-01-26 p
\VUblancAlinea
La tekto dil kavaleyo esas \VUgras{tegulkovrita} \textsuperscript{(55)},\nl
ta dil domi esas \VUgras{ardezkovrita} \textsuperscript{(32)} e ta dil ve\cc
rando esas \VUgras{zinka} \textsuperscript{(24)}.

%%K t05-01-27 p
\VUblancAlinea
\textsc{Ioannes}. --- Ka la tektisto laboras anke la\nl
zinka tekti?

%%K t05-01-28 p
\VUblancAlinea
S\textsuperscript{o} \textsc{Blanko}. --- No, ma la \VUgras{zinko-laboristo} \textsuperscript{(121)}\nl
o la \VUgras{plombo-laboristo}, ta qua pozas \VUgras{luko} \textsuperscript{(38)}\nl
sur la tekto di la domo. Il pozas anke la pluv-\nl
\VUgras{kanali} \textsuperscript{(43)} e la \VUgras{tekto-kanali} \textsuperscript{(42)}. Il kunsoldas\nl
la zinko e la lado. Il fixigis la \VUgras{parafulmino} \textsuperscript{(40)}\nl
e la \VUgras{vento-flecho} \textsuperscript{(39)} sur la \VUgras{geblo} \textsuperscript{(36)}.

%%K t05-01-29 p
\VUblancAlinea
\textsc{Ioannes}. --- Tale la domo esas parfinita?

%%K t05-tit-5 sub
\VUsaut{3.0mm}
\VUpk{10.2pt}{40}{La utensili.}
\VUsaut{3.0mm}

%%K t05-01-30 p
\VUblancAlinea
S\textsuperscript{o} \textsc{Blanko}. --- Ne komplete. La \VUgras{gipsisto} \textsuperscript{(99)}\nl
konstruktas per \VUgras{briki} \textsuperscript{(61)} la parieti qui dividas\nl
olu a diversa chambri. Il indutas per \VUgras{gipso} \textsuperscript{(70)}\nl
la \VUgras{plafoni} \textsuperscript{(26)} e la muri. Nun il laboras la pla\cc
fono di la \VUgras{verando} \textsuperscript{(33)}. Il havas, quale la ma\cc
sonisto, \VUgras{trulo} \textsuperscript{(100)} e \VUgras{trogo} \textsuperscript{(101)}.

%%K t05-01-31 p
\VUblancAlinea
Pose la menuzisto laboras la pordi, la fenes\cc
tri, la \VUgras{parqueti} \textsuperscript{(16)} e la \VUgras{shutri} \textsuperscript{(19)}.

%%K t05-01-32 p
\VUblancAlinea
Nun il laboras en la korto sur sua labor-\nl
\VUgras{benko} \textsuperscript{(90)}. Il havas \VUgras{segilo} \textsuperscript{(94)} por segar la li\cc
gno (69), \VUgras{rabotilo} \textsuperscript{(91)}, \VUgras{fixigilo} \textsuperscript{(92)}, \VUgras{cizelo} \textsuperscript{(94 bis)},\nl
\VUgras{kompaso} \textsuperscript{(95)}, e ligna \VUgras{martelo} \textsuperscript{(95 bis)}.
\end{VUpage}

% --- Feuillet 36 (folio 34) ------------------------------------------
\begin{VUpage}[36]{34}
%%K t05-01-33 p
\VUblancAlinea
\textsc{Ioannes}. --- Ha! ma esas plu amuzanta me\cc
nuzar kam masonar. Onklo mea, tu kompros\nl
por me labor-benko, segilo e rabotilo, nam me\nl
balde laboros kabano por Medor.

%%K t05-01-34 p
\VUblancAlinea
\textsc{L'Onklo}. --- Yes, puereto.

%%K t05-01-35 p
\VUblancAlinea
\textsc{Ioannes}. --- Quante kontenta me esas! E qua\nl
esas ta qua stacas proxim l'eniro-pordo?

%%K t05-tit-6 sub
\VUsaut{3.0mm}
\VUpk{10.2pt}{40}{L'indutado.}
\VUsaut{3.0mm}

%%K t05-01-36 p
\VUblancAlinea
S\textsuperscript{o} \textsc{Blanko}. --- Il esas \VUgras{indutisto} \textsuperscript{(102)}. Il stacas\nl
sur lua skalo kun poto de \VUgras{farbo} \textsuperscript{(103)} e sua \VUgras{pin}\cc
\VUgras{selo} \textsuperscript{(104)}. Il indutas la ligno dil eniro-pordo\nl
por plu longa konserveso.

%%K t05-01-37 p
\VUblancAlinea
Lu anke glutinas la paperi en la aparta\cc
menti.

%%K t05-01-38 p
\VUblancAlinea
La \VUgras{tapetisto} \textsuperscript{(107)} qua esas sur \VUgras{skalo} \textsuperscript{(106)}, sur\nl
la \VUgras{balkono} \textsuperscript{(28)}, pozas \VUgras{storo} \textsuperscript{(29)}.

%%K t05-01-39 p
\VUblancAlinea
La laboristo stacanta proxim la granda fe\cc
nestro esas la \VUgras{vitristo} \textsuperscript{(96)}. Il fixigas la vitro-\nl
\VUgras{kareli} \textsuperscript{(18)} per \VUgras{mastico} \textsuperscript{(73)}. Ta altra laboristo,\nl
genupozinta proxim la pordo di la kelero esas\nl
\VUgras{seruristo} \textsuperscript{(97)} Il pozas \VUgras{seruro} \textsuperscript{(6)}. Lua \VUgras{limilo} \textsuperscript{(98)}\nl
esas proxim ilu.

%%K t05-01-40 p
\VUblancAlinea
\textsc{Ioannes}. --- Kad esas anke seruristo, ta labo\cc
risto qua frapas per sua \VUgras{martelo} \textsuperscript{(84)} sur ferajo?

%%K t05-01-41 p
\VUblancAlinea
S\textsuperscript{o} \textsc{Blanko}. --- Plujustadice, lu esas for\cc
jisto (125). Il forjas \VUgras{feraji} \textsuperscript{(67)} por l'\VUgras{elektristo} \textsuperscript{(124)}\nl
qua esas sur la tekto di la \VUgras{vetureyo} \textsuperscript{(51)}. Il ha\cc
vas \VUgras{forjeyo} \textsuperscript{(126)}, \VUgras{amboso} \textsuperscript{(127)} e \VUgras{pinco} \textsuperscript{(128)}.

%%K t05-01-42 p
\VUblancAlinea
L'elektristo fixigas la \VUgras{kupra} \textsuperscript{(44)} \VUgras{filo} dil tele\cc
fono.
\end{VUpage}

% --- Feuillet 37 (folio 35) ------------------------------------------
\begin{VUpage}[37]{35}
%%K t05-tit-7 sub
\VUpk{10.2pt}{40}{La garden-kultivado.}
\VUsaut{3.0mm}

%%K t05-01-43 p
\VUblancAlinea
\textsc{L'Onklo}. --- Funde dil \VUgras{korto} \textsuperscript{(45)}, proxim la\nl
\VUgras{greto} \textsuperscript{(47)} e la \VUgras{pordego} \textsuperscript{(48)} tu vidas la \VUgras{garde}\cc
\VUgras{nisto} \textsuperscript{(129)}. Il preparas la \VUgras{gardeneto} \textsuperscript{(49)}. Il spa\cc
dagis la tero per sua \VUgras{spado} \textsuperscript{(131)}, il semas se\cc
mini e rastas per la \VUgras{rastilo} \textsuperscript{(130)}. Pose, per sua\nl
\VUgras{arozilo} \textsuperscript{(132)}, il aquizos la \VUgras{bedi} \textsuperscript{(150)} e la planti\nl
kreskados.

%%K t05-01-44 p
\VUblancAlinea
La \VUgras{kaval-servisto} \textsuperscript{(133)} qua jus flegis la ka\cc
valo e donis a lu nutrivo netigas la \VUgras{veturo} \textsuperscript{(52)}\nl
per sponjo, \VUgras{manupumpilo} \textsuperscript{(134)} ed aquo-siteledo.\nl
Pose il rienirigos lu en la vetureyo e klozos la\nl
pordo per la dika \VUgras{riglo} \textsuperscript{(53)}.

%%K t05-01-45 p
\VUblancAlinea
Yen tu kontenta, me opinionas, tu havis\nl
suficant expliki.

%%K t05-01-46 p
\VUblancAlinea
\textsc{Ioannes}. --- Yes, onklo mea, ma volunte me\nl
vidus l'internajo di la domo.

%%K t05-01-47 p
\VUblancAlinea
\textsc{L'Onklo}. --- To eventos morge. Siorulo\nl
Rufo explikos a tu la plano ed indikos a tu\nl
la nomo di singla chambro.

%%K t05-tit-8 sub
\VUsaut{3.0mm}
\VUpk{10.2pt}{40}{Interna aranjo dil domo.}
\VUsaut{2.4mm}

%%K t05-apar-2 apar
{\centering\textit{(Videz la plano.)}\par}
\VUsaut{2.4mm}

%%K t05-01-48 p
\VUblancAlinea
\textsc{L'Arkitekto}. --- Venez, Ioannes, me quik\nl
igos tu explorar la diversa parti di la domo.

%%K t05-01-49 p
\VUblancAlinea
Ni enirez la \VUgras{teretajo} \textsuperscript{(7)}. Yen la butono di\nl
la \VUgras{klosheto} \textsuperscript{(11)}. Ni acensas la tri \VUgras{gradi} \textsuperscript{(14)} dil\nl
\VUgras{perono} \textsuperscript{(9)}, ni transiras la \VUgras{solio} \textsuperscript{(10)} dil eniro-\nl
\VUgras{pordo} \textsuperscript{(8)} e yen ke ni esas an la pedo dil \VUgras{eska}\cc
\VUgras{lero} \textsuperscript{(13)}.

%%K t05-01-50 p
\VUblancAlinea
\textsc{Ioannes}. --- To esas la koridoro.
\end{VUpage}

% --- Feuillet 38 (folio 36) ------------------------------------------
\begin{VUpage}[38]{36}
%%K t05-01-51 p
\VUblancAlinea
\textsc{L'Arkitekto}. --- No, to esas la \VUgras{vestibulo} \textsuperscript{(12)}.\nl
La \VUgras{koridoro} \textsuperscript{(a)} esas ye l'extremajo. Dextre,\nl
yen la \VUgras{salono} \textsuperscript{(b)} e la \VUgras{manjo-chambro} \textsuperscript{(c)}, opoze\nl
a la manjo-chambro esas la \VUgras{koqueyo} \textsuperscript{(d)} e la\nl
\VUgras{servisteyo} \textsuperscript{(e)}, pose avane la \VUgras{skribo-chambro}\textsuperscript{(f)}.\nl
La \VUgras{verando} \textsuperscript{(23)} kun sua \VUgras{vitrizita pordo} \textsuperscript{(25)}\nl
esas proxim la salono.

%%K t05-01-52 p
\VUblancAlinea
Ni quik acensos l'eskalero, Sustenez tu per\nl
la \VUgras{pasamano} \textsuperscript{(15)} e kontez la gradi. Li esas dek-\nl
e-quar. Yen la \VUgras{fluro} \textsuperscript{(m)}, ni esas sur l'unesma\nl
\VUgras{etajo} \textsuperscript{(27)}.

%%K t05-tit-9 sub
\VUsaut{3.0mm}
\VUpk{10.2pt}{40}{L'unesma etajo.}
\VUsaut{3.0mm}

%%K t05-01-53 p
\VUblancAlinea
Opoze a l'eskalero, esas la \VUgras{latrino} \textsuperscript{(20)} e la\nl
\VUgras{tualeteyo} \textsuperscript{(j)}. Dextre bela \VUgras{dormo-chambro} \textsuperscript{(g)}\nl
kun du fenestri an la \VUgras{fasado} \textsuperscript{(1)} di la domo.\nl
Sinistre, esas la \VUgras{balneyo} \textsuperscript{(1)} e la \VUgras{chambro por}\nl
\VUgras{amiki} \textsuperscript{(i)} kun grand \VUgras{alkovo} \textsuperscript{(h)}: Proxim la dor\cc
mo-chambro esas la \VUgras{chambro por l'infanti} \textsuperscript{(k)}.\nl
Ol situesas an vasta \VUgras{teraso} \textsuperscript{(21)}. Sub ica teraso\nl
esas altra latrino (20) e, an la muro, yen la\nl
\VUgras{klosho} \textsuperscript{(22)} qua sonas por la dineo.

%%K t05-01-54 p
\VUblancAlinea
\textsc{Ioannes}. --- Ni irez sur la \VUgras{duesma etajo}.

%%K t05-01-55 p
\VUblancAlinea
\textsc{L'Arkitekto}. --- Ne existas duesma etajo.\nl
Sub la tekto esas nur \VUgras{mansardi} \textsuperscript{(33)} e la gra\cc
nario (34). Ni finis la exploro, ni povas retro\cc
decensar.

%%K t05-01-56 p
\VUblancAlinea
\textsc{Ioannes}. --- Danko, sioro, to esas tre inte\cc
resanta. Me quik naracos a mea patrulo omno\nl
quon me vidis.

\VUsaut{4.0mm}
\VUfilet{22mm}
\end{VUpage}
